% ---- score2 standard preamble (from: python3 score2_tools.py preamble) ----
\documentclass[10pt]{extarticle}
\usepackage[a4paper,margin=1.5cm]{geometry}
\usepackage[T1]{fontenc}
\usepackage{amsmath,amssymb}
\usepackage{graphicx}
\usepackage[export]{adjustbox}
\usepackage{multicol}
\usepackage{enumitem}
\usepackage{booktabs,array,tabularx}
\setlength{\columnsep}{0.7cm}
\setlength{\parindent}{0pt}
\setlength{\parskip}{3pt}
\pagestyle{empty}
\makeatletter
% figure registry (figdefs.tex is generated by the extract step)
\newcommand{\qfigdef}[4]{\@namedef{qf@file@#1}{#2}\@namedef{qf@w@#1}{#3}\@namedef{qf@wide@#1}{#4}}
\newcommand{\qf@check}[1]{\@ifundefined{qf@file@#1}{\PackageError{score2}{Unknown figure id #1}{}}{}}
\newcommand{\qf@img}[4][]{\includegraphics[width=#2\dimexpr\@nameuse{qf@w@#4}pt\relax,max width=#3,max height=0.45\textheight,#1]{\@nameuse{qf@file@#4}}}
% Full-width questions: finalize wraps every question that contains a wide figure in \qfullbegin...\qfullend,
% so the WHOLE question (stem, figures, options) is set across the page in reading order, then the two columns
% resume. The question's head - its number, stem and first figure (\qfullsplit marks the end of the head) -
% is kept on one page, so a number is never left at a page foot with its figure overleaf; the rest may break
% across pages normally, so no big gaps are left. Without a \qfullsplit, the whole question is kept together
% if it is under 45% of a page. Never write these yourself.
\newif\ifqf@full
\newif\ifqf@boxed
\newsavebox{\qf@box}
\newcommand{\qf@open}{\global\qf@fulltrue\par\global\qf@boxedtrue
  \setbox\qf@box\vbox\bgroup\hsize\textwidth\linewidth\textwidth}
\newcommand{\qf@emit}[1]{\par\egroup\global\qf@boxedfalse
  \ifdim\dimexpr\ht\qf@box+\dp\qf@box\relax<#1\textheight\noindent\usebox{\qf@box}\par\else\unvbox\qf@box\fi}
\newcommand{\qfullsplit}{\ifqf@boxed\qf@emit{0.85}\fi}
\newcommand{\qf@close}{\ifqf@boxed\qf@emit{0.45}\fi\global\qf@fullfalse}
\newcommand{\qfullbegin}{\end{multicols}\qf@open}
\newcommand{\qfullend}{\qf@close\begin{multicols}{2}}
\newcommand{\qfullstart}{\qf@open}
\newcommand{\qfulldone}{\qf@close}
% \qfig[scale]{id}: block figure, centred, directly where it is written (size set by the tool).
% Fallback if a wide figure is not inside a full-width question: break out of the columns in place - no floats.
\newcommand{\qfig}[2][1]{\qf@check{#2}%
  \ifqf@full\qf@blk{#1}{#2}\else
  \ifnum\@nameuse{qf@wide@#2}=1 \expandafter\@firstoftwo\else\expandafter\@secondoftwo\fi
  {\end{multicols}\begin{center}\qf@img{#1}{\textwidth}{#2}\end{center}\begin{multicols}{2}}%
  {\qf@blk{#1}{#2}}\fi}
\newcommand{\qf@blk}[2]{\par\smallskip{\centering\qf@img{#1}{\linewidth}{#2}\par}\smallskip}
% \qfigi[scale]{id}: inline figure (inside option lists / tables / minipages); never breaks columns.
\newcommand{\qfigi}[2][1]{\qf@check{#2}\qf@img[valign=c]{#1}{\linewidth}{#2}}
\makeatother
% structure macros - lint reads these to check completeness and figure placement
\newcommand{\papersec}[1]{\par\bigskip{\centering\Large\bfseries PAPER-#1\par}\medskip}
\newcommand{\partsec}[2]{\par\medskip\noindent{\large\bfseries PART-#1: #2}\par\smallskip}
\newcommand{\sect}[1]{\par\smallskip\noindent\textbf{#1}\par}
\newcommand{\qn}[1]{\par\smallskip\noindent\textbf{#1.}~}
\newcommand{\paraq}[2]{\par\smallskip\noindent\textbf{Paragraph for Questions #1 and #2}\par}
\newenvironment{qoptions}{\begin{itemize}[leftmargin=1.8em,itemsep=1pt,topsep=2pt,label={}]}{\end{itemize}}
\newcommand{\opt}[2]{\item[(#1)] #2}
% figure registry (generated by score2_tools.py finalize)
\qfigdef{p1_phy_s1_q02_1}{Score-II - Test 39 (02 Mar 2025)-img/p1_phy_s1_q02_1.png}{103.8}{0}
\qfigdef{p1_phy_s1_q03_1}{Score-II - Test 39 (02 Mar 2025)-img/p1_phy_s1_q03_1.png}{100.4}{0}
\qfigdef{p1_phy_s1_q05_1}{Score-II - Test 39 (02 Mar 2025)-img/p1_phy_s1_q05_1.png}{109.3}{0}
\qfigdef{p1_phy_s2_q08_1}{Score-II - Test 39 (02 Mar 2025)-img/p1_phy_s2_q08_1.png}{106.9}{0}
\qfigdef{p1_phy_s2_q08_2}{Score-II - Test 39 (02 Mar 2025)-img/p1_phy_s2_q08_2.png}{106.9}{0}
\qfigdef{p1_phy_s2_q08_3}{Score-II - Test 39 (02 Mar 2025)-img/p1_phy_s2_q08_3.png}{106.9}{0}
\qfigdef{p1_phy_s2_q08_4}{Score-II - Test 39 (02 Mar 2025)-img/p1_phy_s2_q08_4.png}{107.1}{0}
\qfigdef{p1_phy_s2_q09_1}{Score-II - Test 39 (02 Mar 2025)-img/p1_phy_s2_q09_1.png}{108.9}{0}
\qfigdef{p1_phy_s2_q10_1}{Score-II - Test 39 (02 Mar 2025)-img/p1_phy_s2_q10_1.png}{89.1}{0}
\qfigdef{p1_phy_s3_q01_1}{Score-II - Test 39 (02 Mar 2025)-img/p1_phy_s3_q01_1.png}{109.6}{0}
\qfigdef{p1_phy_s3_q02_1}{Score-II - Test 39 (02 Mar 2025)-img/p1_phy_s3_q02_1.png}{89.0}{0}
\qfigdef{p1_phy_s3_q05_1}{Score-II - Test 39 (02 Mar 2025)-img/p1_phy_s3_q05_1.png}{143.3}{0}
\qfigdef{p1_phy_s3_q06_1}{Score-II - Test 39 (02 Mar 2025)-img/p1_phy_s3_q06_1.png}{119.6}{0}
\qfigdef{p1_phy_s3_q07_1}{Score-II - Test 39 (02 Mar 2025)-img/p1_phy_s3_q07_1.png}{145.6}{0}
\qfigdef{p1_chem_s1_q05_1}{Score-II - Test 39 (02 Mar 2025)-img/p1_chem_s1_q05_1.png}{54.9}{0}
\qfigdef{p1_chem_s1_q05_2}{Score-II - Test 39 (02 Mar 2025)-img/p1_chem_s1_q05_2.png}{49.2}{0}
\qfigdef{p1_chem_s1_q05_B}{Score-II - Test 39 (02 Mar 2025)-img/p1_chem_s1_q05_B.png}{53.7}{0}
\qfigdef{p1_chem_s1_q05_C}{Score-II - Test 39 (02 Mar 2025)-img/p1_chem_s1_q05_C.png}{53.6}{0}
\qfigdef{p1_chem_s1_q05_D}{Score-II - Test 39 (02 Mar 2025)-img/p1_chem_s1_q05_D.png}{46.8}{0}
\qfigdef{p1_chem_s1_q05_6}{Score-II - Test 39 (02 Mar 2025)-img/p1_chem_s1_q05_6.png}{230.4}{0}
\qfigdef{p1_chem_s3_q07_1}{Score-II - Test 39 (02 Mar 2025)-img/p1_chem_s3_q07_1.png}{325.4}{1}
\qfigdef{p2_phy_s1_q01_1}{Score-II - Test 39 (02 Mar 2025)-img/p2_phy_s1_q01_1.png}{82.6}{0}
\qfigdef{p2_phy_s1_q01_2}{Score-II - Test 39 (02 Mar 2025)-img/p2_phy_s1_q01_2.png}{88.6}{0}
\qfigdef{p2_phy_s1_q01_B}{Score-II - Test 39 (02 Mar 2025)-img/p2_phy_s1_q01_B.png}{90.4}{0}
\qfigdef{p2_phy_s1_q01_C}{Score-II - Test 39 (02 Mar 2025)-img/p2_phy_s1_q01_C.png}{89.1}{0}
\qfigdef{p2_phy_s1_q01_D}{Score-II - Test 39 (02 Mar 2025)-img/p2_phy_s1_q01_D.png}{91.0}{0}
\qfigdef{p2_phy_s2_q06_1}{Score-II - Test 39 (02 Mar 2025)-img/p2_phy_s2_q06_1.png}{143.1}{0}
\qfigdef{p2_phy_s2_q07_1}{Score-II - Test 39 (02 Mar 2025)-img/p2_phy_s2_q07_1.png}{36.9}{0}
\qfigdef{p2_phy_s2_q08_1}{Score-II - Test 39 (02 Mar 2025)-img/p2_phy_s2_q08_1.png}{117.9}{0}
\qfigdef{p2_phy_s2_q09_1}{Score-II - Test 39 (02 Mar 2025)-img/p2_phy_s2_q09_1.png}{77.1}{0}
\qfigdef{p2_phy_s2_q10_1}{Score-II - Test 39 (02 Mar 2025)-img/p2_phy_s2_q10_1.png}{104.1}{0}
\qfigdef{p2_phy_s3_q01_1}{Score-II - Test 39 (02 Mar 2025)-img/p2_phy_s3_q01_1.png}{95.4}{0}
\qfigdef{p2_phy_s3_q02_1}{Score-II - Test 39 (02 Mar 2025)-img/p2_phy_s3_q02_1.png}{89.3}{0}
\qfigdef{p2_phy_s3_q03_1}{Score-II - Test 39 (02 Mar 2025)-img/p2_phy_s3_q03_1.png}{105.7}{0}
\qfigdef{p2_phy_s3_q05_1}{Score-II - Test 39 (02 Mar 2025)-img/p2_phy_s3_q05_1.png}{108.5}{0}
\qfigdef{p2_phy_s3_q08_1}{Score-II - Test 39 (02 Mar 2025)-img/p2_phy_s3_q08_1.png}{134.9}{0}
\qfigdef{p2_chem_s1_q04_1}{Score-II - Test 39 (02 Mar 2025)-img/p2_chem_s1_q04_1.png}{280.1}{1}
\qfigdef{p2_chem_s1_q04_3}{Score-II - Test 39 (02 Mar 2025)-img/p2_chem_s1_q04_3.png}{38.1}{0}
\qfigdef{p2_chem_s1_q04_B}{Score-II - Test 39 (02 Mar 2025)-img/p2_chem_s1_q04_B.png}{48.9}{0}
\qfigdef{p2_chem_s1_q04_C}{Score-II - Test 39 (02 Mar 2025)-img/p2_chem_s1_q04_C.png}{63.0}{0}
\qfigdef{p2_chem_s1_q04_D}{Score-II - Test 39 (02 Mar 2025)-img/p2_chem_s1_q04_D.png}{49.8}{0}
\qfigdef{p2_chem_s2_q05_1}{Score-II - Test 39 (02 Mar 2025)-img/p2_chem_s2_q05_1.png}{109.0}{0}
\qfigdef{p2_chem_s2_q09_1}{Score-II - Test 39 (02 Mar 2025)-img/p2_chem_s2_q09_1.png}{199.0}{0}
\qfigdef{p2_chem_s2_q10_1}{Score-II - Test 39 (02 Mar 2025)-img/p2_chem_s2_q10_1.png}{67.1}{0}
\qfigdef{p2_chem_s2_q10_B}{Score-II - Test 39 (02 Mar 2025)-img/p2_chem_s2_q10_B.png}{173.1}{0}
\qfigdef{p2_chem_s2_q10_C}{Score-II - Test 39 (02 Mar 2025)-img/p2_chem_s2_q10_C.png}{302.1}{1}
\qfigdef{p2_chem_s2_q10_D}{Score-II - Test 39 (02 Mar 2025)-img/p2_chem_s2_q10_D.png}{174.0}{0}
\qfigdef{p2_chem_s3_q06_1}{Score-II - Test 39 (02 Mar 2025)-img/p2_chem_s3_q06_1.png}{245.2}{0}
\qfigdef{p2_chem_s3_q07_1}{Score-II - Test 39 (02 Mar 2025)-img/p2_chem_s3_q07_1.png}{115.7}{0}
\qfigdef{p2_chem_s3_q08_1}{Score-II - Test 39 (02 Mar 2025)-img/p2_chem_s3_q08_1.png}{197.3}{0}
\raggedcolumns
\makeatletter\@beginparpenalty=10000\makeatother
\begin{document}
\begin{center}{\Large\bfseries JEE (Main+Advanced) Score-II --- Test 39 (02 Mar 2025)}\end{center}
\begin{multicols}{2}

\papersec{1}
\partsec{1}{PHYSICS}
\sect{Section-I(i)}

\qn{1} A conducting spherical shell of inner radius a and outer radius b is made up of a material of resistivity $\rho$ and negligible dielectric activity. At time t = 0, material of shell is electrically neutral and a charge $q_0$ is some how located at the centre of shell (i.e. at time = 0, there are no induced charges, they will appear afterward).\\
Choose correct statement(s) :
\begin{qoptions}
\opt{A}{When $-Q_t$ charge is induced at inner surface at time t, then electric field in the material of shell is $E=\dfrac{q_0-Q_t}{4\pi\epsilon_0r^2}$ where r is distance of point under consideration from centre of shell.}
\opt{B}{Current through shell at this moment (discussed in option A) is $I=\dfrac{q_0-Q_t}{\epsilon_0\rho}$}
\opt{C}{At time t we can write $\dfrac{dQ_t}{q_0-Q_t}=\dfrac{dt}{\epsilon_0\rho}$}
\opt{D}{At time t, value of $Q_t$ is given by $Q_t=q_0\left(1-e^{-\frac{t}{\epsilon_0\rho}}\right)$}
\end{qoptions}

\qn{2} A bottle is kept on the ground as shown in the figure. The bottle can be modelled as having two cylindrical zones. The lower zone of the bottle has a cross-sectional radius of $R\sqrt2$ and is filled with honey of density $2\rho$. The upper zone of the bottle is filled with the water of density $\rho$ and has a cross-sectional radius R. The height of the lower zone is H while that of the upper zone is 2H. If now the honey and the water parts are mixed together to form a homogeneous solution. (Assume that total volume does not change)
\qfig{p1_phy_s1_q02_1}
\begin{qoptions}
\opt{A}{The pressure inside the bottle at the base will remains unaltered}
\opt{B}{The normal reaction on the bottle from the ground will remain unaltered}
\opt{C}{The pressure inside the bottle at the base will increase by an amount $\left(\dfrac12\right)\rho gH$}
\opt{D}{The pressure inside the bottle at the base will decrease by an amount $\left(\dfrac14\right)\rho gH$}
\end{qoptions}

\qn{3} In the given $P$-$V$ diagram, a monoatomic gas $\left(\gamma=\dfrac53\right)$ is first compressed adiabatically from state $A$ to state $B$. Then it expands isothermally from state $B$ to state $C$. [Given: $\left(\dfrac13\right)^{0.6}\simeq0.5$, $\ln2\simeq0.7$]. Which of the following statement(s) is(are) correct?
\qfig{p1_phy_s1_q03_1}
\begin{qoptions}
\opt{A}{The magnitude of the total work done in the process $A\to B\to C$ is 144 kJ.}
\opt{B}{The magnitude of the work done in the process $B\to C$ is 84 kJ.}
\opt{C}{The magnitude of the work done in the process $A\to B$ is 60 kJ.}
\opt{D}{The magnitude of the work done in the process $C\to A$ is zero.}
\end{qoptions}

\qn{4} Choose the correct statement(s) :-
\begin{qoptions}
\opt{A}{An ellipse is drawn in equatorial plane of a dipole. All the points on ellipse, will be equipotential.}
\opt{B}{Electric field at a distance `\textbf{r}' from small dipole is 5 N/C at a point. Any other point at a distance `\textbf{r}' from dipole can have electric field strength equal to 15 N/C.}
\opt{C}{Force on a dipole in external electric field will be zero always.}
\opt{D}{Torque on a dipole in external electric field will be non zero.}
\end{qoptions}

\qn{5} A uniform thin plate of density $\rho$ kg/m$^3$ is shown in the figure. Then which of the following statements is/are correct?
\qfig{p1_phy_s1_q05_1}
\begin{qoptions}
\opt{A}{Moment of inertia of plate about x-axis is $\dfrac{16\rho a^7}{21}$}
\opt{B}{Moment of inertia of the plate about y-axis is $\dfrac{2\rho a^5}5$}
\opt{C}{Moment of inertia of the plate about y-axis is $\dfrac{4\rho a^5}5$}
\opt{D}{Moment of inertia of the plate about x-axis is $\dfrac{8\rho a^7}7$}
\end{qoptions}

\qn{6} A particle of mass m and charge q is projected into a region having a uniform magnetic field $B_0$. Initial velocity ($v_0$) of the particle is perpendicular to the magnetic field. Apart from the magnetic force, the particle faces a frictional force which has a magnitude of f = kv where v is instantaneous speed and k is a positive constant. Considering initial time of motion at t = 0. Select the correct option(s):
\begin{qoptions}
\opt{A}{Particle will follow a circular path of constant radius.}
\opt{B}{Radius of curvature of the path of the particle at time `t' will be $\dfrac{mv_0}{qB}e^{-\frac{kt}{2m}}$}
\opt{C}{Radius of curvature of the path of the particle at time `t' will be $\dfrac{mv_0}{qB}e^{-\frac{kt}{m}}$}
\opt{D}{After travelling through a distance of $x=\dfrac{mv_0}{2k}$, radius of curvature of the path of the particle will be $R=\dfrac{mv_0}{2qB}$.}
\end{qoptions}

\sect{Section-I(ii)}

\qn{7} A charged particle having non-zero velocity is subjected to certain conditions given in List-I. List-II gives possible trajectories of the particle. Match the conditions in List-I with results in List-II\par\nobreak
\begin{center}\small\setlength{\tabcolsep}{3pt}
\begin{tabular}{|c|p{3.0cm}|c|p{2.8cm}|}
\hline
 & \multicolumn{1}{c|}{\textbf{List-I}} & & \multicolumn{1}{c|}{\textbf{List-II}} \\
\hline
(I) & In only uniform electric field & (P) & The path may be a straight line \\
\hline
(II) & In only uniform magnetic field & (Q) & The path may be a parabola \\
\hline
(III) & In uniform magnetic and uniform electric field that are both parallel & (R) & The path may be a circle \\
\hline
(IV) & Subjected to a net force of constant magnitude & (S) & The path may be a helix with uniform pitch \\
\hline
 & & (T) & The path may be a helix with non-uniform pitch \\
\hline
\end{tabular}
\end{center}
\begin{qoptions}
\opt{A}{I $\to$ P,Q; II $\to$ P,R; III $\to$ P,R,T; IV $\to$ P,Q,R,S,T}
\opt{B}{I $\to$ P,Q; II $\to$ P,R; III $\to$ P,Q,R; IV $\to$ P}
\opt{C}{I $\to$ P,Q; II $\to$ P,R,S; III $\to$ P,T; IV $\to$ P,Q,R,S,T}
\opt{D}{I $\to$ P; II $\to$ P,Q; III $\to$ P,S,T; IV $\to$ P,Q,R,S,T}
\end{qoptions}

\qn{8} In all the situations shown in list-I blocks are connected with ideal spring of constant K. Initially they are at rest and spring is unstretched (Neglect friction). Match the corresponding options in list-II.\par\nobreak
\begin{center}\small\setlength{\tabcolsep}{3pt}
\begin{tabular}{|c|p{3.3cm}|c|p{2.4cm}|}
\hline
 & \multicolumn{1}{c|}{\textbf{List-I}} & & \multicolumn{1}{c|}{\textbf{List-II}} \\
\hline
(I) & \qfigi{p1_phy_s2_q08_1} & (P) & Maximum elongation in the spring is 4F/3K. \\
\hline
(II) & \qfigi{p1_phy_s2_q08_2} & (Q) & Magnitude of acceleration of centre of mass of two block system is $\dfrac F{3m}$ \\
\hline
(III) & \qfigi{p1_phy_s2_q08_3} & (R) & Relative velocity of block is maximum when elongation in the spring $\dfrac FK$ \\
\hline
(IV) & \qfigi{p1_phy_s2_q08_4} & (S) & Maximum magnitude of relative velocity between the blocks is $\dfrac Fm\sqrt{\dfrac{6m}K}$ \\
\hline
 & & (T) & None of these \\
\hline
\end{tabular}
\end{center}
\begin{qoptions}
\opt{A}{I $\to$ P; II $\to$ Q; III $\to$ R; IV $\to$ S}
\opt{B}{I $\to$ Q; II $\to$ S; III $\to$ P; IV $\to$ R}
\opt{C}{I $\to$ R; II $\to$ Q; III $\to$ S; IV $\to$ P}
\opt{D}{I $\to$ S; II $\to$ P; III $\to$ Q; IV $\to$ R}
\end{qoptions}

\qn{9} A tall cylinder partially filled with mercury (Hg) is rotated about its axis with a constant angular velocity $\omega_0$ as shown in the figure. If $\rho(x)$ is the radius of curvature at any co-ordinate x (assuming lowest point of Hg surface at origin) and h(x) is the height of free surface of Hg measured at a coordinate `x' above the `x' axis. Then match List-I with List-II.
\qfig{p1_phy_s2_q09_1}\par\nobreak
\begin{center}\small\setlength{\tabcolsep}{3pt}
\begin{tabular}{|c|p{1.4cm}|c|p{2.9cm}|}
\hline
 & \multicolumn{1}{c|}{\textbf{List-I}} & & \multicolumn{1}{c|}{\textbf{List-II}} \\
\hline
(I) & $h\left(\dfrac{r_0}2\right)$ & (P) & $\dfrac{\omega_0^2.r_0^2}{8g}$ \\
\hline
(II) & $h(r_0)$ & (Q) & $\dfrac{\left(1+\frac{\omega_0^4r_0^2}{4g^2}\right)^{3/2}}{\frac{\omega_0^2}g}$ \\
\hline
(III) & $\rho\left(\dfrac{r_0}2\right)$ & (R) & $\dfrac{\omega_0^2r_0^2}{2g}$ \\
\hline
(IV) & $\rho(0)$ & (S) & $\dfrac g{\omega_0^2}$ \\
\hline
 & & (T) & $\dfrac{2g}{\omega_0^2}$ \\
\hline
\end{tabular}
\end{center}
\begin{qoptions}
\opt{A}{I $\to$ R; II $\to$ S; III $\to$ Q; IV $\to$ T}
\opt{B}{I $\to$ Q; II $\to$ R; III $\to$ T; IV $\to$ P}
\opt{C}{I $\to$ P; II $\to$ R; III $\to$ Q; IV $\to$ S}
\opt{D}{I $\to$ S; II $\to$ T; III $\to$ P; IV $\to$ S}
\end{qoptions}

\qn{10} Two parallel plates of width `W' (see figure) are vertically inserted in a perfectly wetting liquid of density `$\rho$'. The separation between the plates is `d' (d $\ll$ W). The liquid rises in between the plates and becomes static. If surfaces tension of liquid is `S' then
\qfig{p1_phy_s2_q10_1}\par\nobreak
\begin{center}\small\setlength{\tabcolsep}{3pt}
\begin{tabular}{|c|p{2.8cm}|c|p{1.5cm}|}
\hline
 & \multicolumn{1}{c|}{\textbf{List-I}} & & \multicolumn{1}{c|}{\textbf{List-II}} \\
\hline
(I) & Height risen in between the plates is given by & (P) & Zero \\
\hline
(II) & The heat energy released in the process of rising liquid & (Q) & $\dfrac{4S^2\omega}{\rho gd}$ \\
\hline
(III) & The force of attraction between the plates & (R) & $\dfrac{2S}{\rho gd}$ \\
\hline
(IV) & The work done by surface tension in raising the liquid & (S) & $\dfrac{2S^2\omega}{\rho gd^2}$ \\
\hline
 & & (T) & $\dfrac{2S^2\omega}{\rho gd}$ \\
\hline
\end{tabular}
\end{center}
\begin{qoptions}
\opt{A}{I $\to$ R; II $\to$ S; III $\to$ Q; IV $\to$ P}
\opt{B}{I $\to$ R; II $\to$ P; III $\to$ T; IV $\to$ S}
\opt{C}{I $\to$ R; II $\to$ T; III $\to$ S; IV $\to$ Q}
\opt{D}{I $\to$ R; II $\to$ P; III $\to$ P; IV $\to$ Q}
\end{qoptions}

\sect{Section-II}

\qn{1} A satellite is moving on a circular orbit around a planet. A sharp impulse is given to the satellite such that it just escapes from the gravitational field of the planet. The impulse can be given in two different ways as shown in the figure. The ratio of the impulse given in the two different cases $\left(\dfrac{P_1}{P_2}\right)$ is :-
\qfig{p1_phy_s3_q01_1}

\qn{2} A plank of length 20 m and mass 1 kg is kept on a horizontal smooth surface. A cylinder of mass 1 kg is kept near one end of the plank. The coefficient of friction between the two surfaces is 0.5. Plank is suddenly given a velocity 20 m/s towards left. Find the time (in sec) at which plank and cylinder separate.
\qfig{p1_phy_s3_q02_1}

\qn{3} Shin Chan and his mother have a tin whistle each. The pipe length of Shin chan's tin whistle is 52 cm long while the pipe length of mother's tin whistle is 50 cm long. They both play at the same time, sounding the whistles at their fundamental resonant frequencies. The velocity of sound in air is 325 m/s. Assume the whistle is a pipe with one end open find the beat frequency (in Hz) that is heard when both whistles are playing simultaneously.

\qn{4} A thin film of a specific material can be used to decrease the intensity of reflected light. There is destructive interference of waves reflected from upper and lower surfaces of the film. These films are called anti-reflection coating. Magnesium fluoride (MgF$_2$) is used as anti-reflecting coating on plane glass surface for light having wavelength 500 nm and normal incidence. The minimum thickness of film required is 100 nm. Find refractive index of magnesium fluoride. It is known that refractive index of glass is more than refractive index of magnesium fluoride.

\qn{5} From a thin convex lens of focal length F = 20 cm and aperture diameter of 10 mm a small piece was cut as shown in the figure. The piece is then placed in contact with the remaining part of the lens on the original principal axis of the lens. If the optical centre of the original lens is taken to be origin of co-ordinates and a point object is placed at $(-30,\,0)$, then the image coordinate becomes (x, y). Find the value of x + y\\
(both `x' and `y' are in the lowest fractional form)
\qfig{p1_phy_s3_q05_1}

\qn{6} The circuit is closed for long time so that the capacitors are fully charged. Switch is opened at a certain instant. The heat developed in resistor $R_2$ after that instant is K $\times$ $10^{-4}$ J. Then find the value of K.
\qfig{p1_phy_s3_q06_1}

\qn{7} Three alternating voltage sources $V_1=3\sin\omega t$ volt, $V_2=5\sin(\omega t+\phi_1)$ volt and $V_3=5\sin(\omega t-\phi_2)$ volt connected across a resistance $R=\sqrt{\dfrac73}\ \Omega$ as shown in the figure (where $\phi_1$ and $\phi_2$ corresponds to $30^\circ$ and $127^\circ$ respectively). Find the peak current (in Amp) through the resistor.
\qfig{p1_phy_s3_q07_1}

\qn{8} In an X-ray tube the accelerating voltage is 20 kV. A target is used in the X-ray tube. The wavelength of the L$\alpha$ line is 124 pm. The wavelength of K$\alpha$ line is 54 pm. The energy of the `target' ion with vacancy in `M' shell is 5.5 keV higher than the atom of target. Find the minimum wavelength (in \AA) of the characteristic X-ray that will be emitted by target. [Take hc = 12400 eV\AA]

\partsec{2}{CHEMISTRY}
\sect{Section-I(i)}

\qn{1} The correct options are :
\begin{qoptions}
\opt{A}{Geometry of $[\text{AuCl}_4]^-$ - square planar}
\opt{B}{Hybridisation of metal in $[\text{Fe(H}_2\text{O)}_5\text{NO}]^{+2}$ - sp$^3$d$^2$}
\opt{C}{IUPAC name of $[\text{Co(NH}_3)_4\text{Cl(ONO)}]\text{Cl}$ - Tetraamminochloridonitrito-O-cobalt (III) chloride}
\opt{D}{K$_3$[VF$_6$] - Colourless compound}
\end{qoptions}

\qn{2} Which of the following is/are correct ?
\begin{qoptions}
\opt{A}{Egg albumin is a lyophilic colloidal solution}
\opt{B}{Chemisorption is an endothermic process}
\opt{C}{Gold sol is a negatively charged colloid}
\opt{D}{Charcoal adsorbs more hydrogen gas than chlorine gas.}
\end{qoptions}

\qn{3} Which of the following statement(s) is/are correct?
\begin{qoptions}
\opt{A}{NaCl is a `AB' crystal lattice that can be interpreted to be made up of two individual FCC unit cells of A$^+$ and B$^-$ fused together in such a manner that the corner of one unit cell becomes the edge centre of the other.}
\opt{B}{In a face centred cubic unit cell, the body centre is an octahedral void.}
\opt{C}{In FCC unit cell, octahedral and tetrahedral voids are equal in number.}
\opt{D}{Tetrahedral voids = 2 $\times$ octahedral voids, is valid for both CCP and HCP lattices.}
\end{qoptions}

\qn{4} The correct statement among the following are :-
\begin{qoptions}
\opt{A}{B$_2$H$_6$ can survive in air.}
\opt{B}{BCl$_3$ is better lewis acid as compared to BF$_3$ as well as AlCl$_3$}
\opt{C}{Both Graphite \& inorganic graphite are good conductor of electricity.}
\opt{D}{Al$_2$(CH$_3$)$_6$ has 3c - 2e bond(s)}
\end{qoptions}

\qn{5} On reaction with NaNO$_2$ and HCl, which of the following amino alcohols will yield \textbf{P} ?
\qfig{p1_chem_s1_q05_1}
\begin{qoptions}
\opt{A}{\qfigi{p1_chem_s1_q05_2}}
\opt{B}{\qfigi{p1_chem_s1_q05_B}}
\opt{C}{\qfigi{p1_chem_s1_q05_C}}
\opt{D}{\qfigi{p1_chem_s1_q05_D}}
\end{qoptions}

\qn{6}
\qfig{p1_chem_s1_q05_6}
Which of the following statement(s) are true regarding above reactions ?
\begin{qoptions}
\opt{A}{Rate determining step of all three reactions follow second order kinetics}
\opt{B}{Rate determining step of all three reactions follow first order kinetiecs.}
\opt{C}{In reaction III, rate of reaction is independent on the initial concentration of Br$_2$}
\opt{D}{In reaction III, rate of reaction is dependent on the initial concentration of Br$_2$}
\end{qoptions}

\sect{Section-I(ii)}

\qn{7} Consider the following Lists :\par\nobreak
\begin{center}\small\setlength{\tabcolsep}{3pt}
\begin{tabular}{|c|p{4.2cm}|c|p{0.8cm}|}
\hline
 & \multicolumn{1}{c|}{\textbf{List-I}} & & \textbf{List-II} \\
\hline
(A) & If activation energy is 4.606 kCal, then how many time reaction becomes faster at 727$^\circ$C than at 227$^\circ$C & (P) & 10 \\
\hline
(B) & Rate constant of a reaction is 0.0693 s$^{-1}$, then after how much time in minutes, will be required to change in concentration of reactant from 16 M to 2M & (Q) & 2 \\
\hline
(C) & Order of reaction whose half life is inversely proportional to its initial conc. & (R) & 30 \\
\hline
(D) & Order of reaction of hydrolysis of esters in alkaline medium & (S) & 0.5 \\
\hline
 & & (T) & 4 \\
\hline
\end{tabular}
\end{center}
The correct match is :
\begin{qoptions}
\opt{A}{A $\to$ T; B $\to$ R; C $\to$ Q; D $\to$ Q}
\opt{B}{A $\to$ R; B $\to$ P; C $\to$ S; D $\to$ Q}
\opt{C}{A $\to$ P; B $\to$ S; C $\to$ Q; D $\to$ Q}
\opt{D}{A $\to$ P; B $\to$ R; C $\to$ Q; D $\to$ T}
\end{qoptions}

\qn{8} Match the following\par\nobreak
\begin{center}\small\setlength{\tabcolsep}{3pt}
\begin{tabular}{|c|p{2.6cm}|c|p{2.6cm}|}
\hline
 & \textbf{List-I}\newline (Reaction) & & \textbf{List-II}\newline (Shape of metal containing compound) \\
\hline
(P) & Co$^{+2}$ + HCl (conc.) $\rightarrow$ & (1) & Octahedral (d$^2$sp$^3$) \\
\hline
(Q) & AgCl + KCN(excess) $\rightarrow$ & (2) & Tetrahedral \\
\hline
(R) & Ni$^{+2}$ + DMG$^{-1}$ $\rightarrow$ & (3) & Linear \\
\hline
(S) & Co(OH)$_2$ + NH$_3$(excess) $\rightarrow$ & (4) & square planar \\
\hline
 & & (5) & Octahedral (sp$^3$d$^2$) \\
\hline
\end{tabular}
\end{center}
\begin{qoptions}
\opt{A}{P $\to$ 2; Q $\to$ 4; R $\to$ 5; S $\to$ 1}
\opt{B}{P $\to$ 3; Q $\to$ 1; R $\to$ 3; S $\to$ 5}
\opt{C}{P $\to$ 5; Q $\to$ 2; R $\to$ 1; S $\to$ 3}
\opt{D}{P $\to$ 2; Q $\to$ 3; R $\to$ 4; S $\to$ 5}
\end{qoptions}

\qn{9} Match the following\par\nobreak
\begin{center}\small\setlength{\tabcolsep}{3pt}
\begin{tabular}{|c|p{2.0cm}|c|p{3.3cm}|}
\hline
 & \multicolumn{1}{c|}{\textbf{List-I} (Compounds)} & & \multicolumn{1}{c|}{\textbf{List-II} (Uses)} \\
\hline
(A) & BrF$_3$ & (P) & As a flux \\
\hline
(B) & HCl & (Q) & As an antiseptic \\
\hline
(C) & SO$_2$ & (R) & In enrichment of U$^{235}$ \\
\hline
(D) & Na$_2$B$_4$O$_7$ & (S) & In refining of petroleum \\
\hline
 & & (T) & In extracting glue from bones \\
\hline
\end{tabular}
\end{center}
\begin{qoptions}
\opt{A}{A $\to$ P; B $\to$ R; C $\to$ Q; D $\to$ S}
\opt{B}{A $\to$ T; B $\to$ Q; C $\to$ P; D $\to$ R}
\opt{C}{A $\to$ R; B $\to$ T; C $\to$ S; D $\to$ P}
\opt{D}{A $\to$ Q; B $\to$ P; C $\to$ S; D $\to$ T}
\end{qoptions}

\qn{10} Match the following\par\nobreak
\begin{center}\small\setlength{\tabcolsep}{3pt}
\begin{tabular}{|c|p{1.9cm}|c|p{4.0cm}|}
\hline
\multicolumn{2}{|c|}{\textbf{List-I}} & \multicolumn{2}{c|}{\textbf{List-II}} \\
\hline
(P) & Phenol & (1) & White precipitate obtained with Bromine water \\
\hline
(Q) & Aniline & (2) & Sodium fusion extract of the compound on treatment with FeCl$_3$ gives blood red color \\
\hline
(R) & Methionine & (3) & Addition of compound to a saturated solution of NaHCO$_3$ gives effervescence \\
\hline
(S) & Saccharin & (4) & Sodium fusion extract of the compound on boiling with FeSO$_4$ followed by acidification with cone H$_2$SO$_4$ gives prussian blue color \\
\hline
 & & (5) & Violet colored solution is obtained on treating the compound with neutral FeCl$_3$ \\
\hline
\end{tabular}
\end{center}
\begin{qoptions}
\opt{A}{P $\to$ 1,5; Q $\to$ 1,4; R $\to$ 2,3; S $\to$ 2}
\opt{B}{P $\to$ 3,5; Q $\to$ 4; R $\to$ 2,4; S $\to$ 1,5}
\opt{C}{P $\to$ 3,4,5; Q $\to$ 1,2; R $\to$ 2,3,5; S $\to$ 3,5}
\opt{D}{P $\to$ 5; Q $\to$ 1,2; R $\to$ 2; S $\to$ 2,5}
\end{qoptions}

\sect{Section-II}

\qn{1} The vapour pressure of two miscible liquid A and B are 300 mm and 500 mm Hg respectively. In a flask, 10 mole of A is mixed with 20 moles of B. When B is added in A, A starts polymerizing into a completely insoluble solid by following first order kinetics. After 200 min, 1 mole of a nonvolatile solute is dissolved which arrest the process of polymerization completely. Final vapour pressure of solution is 400 mm of Hg, then what will be the value of rate constant in unit hr$^{-1}$ in polymerization process.\\
\relax[Given ln 2 = 0.70, ln 5 = 1.6]

\qn{2} Calculate the pH of solution prepared by mixing 40 ml of 0.1M acetic acid with 20 ml of 0.1M sodium hydroxide. Given $K_{a(\text{ACETIC ACID})}=10^{-5}$

\qn{3} Solution of 0.05 mole of sugar in 500 gm of water shows $0.279^\circ$C depression in freezing point. Cane sugar solution has inversion as follows.\\
Cane sugar + H$_2$O($\ell$) $\rightarrow$ Glucose + Fructose\\
If the final mass (in gm) of cane sugar left in the solution is X gram, then find the value of (X $-$ 3.55) ?\\
\relax[Given : $K_f$ (H$_2$O) = 1.86 K-kg mol$^{-1}$]\\
Ignore decrease in mass of water due to reaction.

\qn{4} When an electron present in ground state of H-atom absorbs a photon of energy 12.09 eV, then how many times will its de-Broglie wavelength increase ?

\qn{5} The total number of enols possible for the formula C$_4$H$_8$O are

\qn{6} A carbonyl compound of formula C$_9$H$_{10}$O(A), which is a benzene derivative gives orange precipitate with 2,4 -DNP and also gives yellow precipitate with I$_2$ in presence of aqueous NaOH. The total no. of isomers possible for `A' are

\qfullbegin
\qn{7} In how many of the following reactions aldehyde is formed as product?
\qfig{p1_chem_s3_q07_1}
\qfullend

\qn{8} In the scheme given below, X and Y, respectively, are\\
$\text{Metal halide}\xrightarrow{\text{aq. NaOH}}\text{Green precipitate (P)}+\text{Filtrate (Q)}$\\
$\text{P}\xrightarrow{\text{Na}_2\text{O}_2\text{ (aq)}}\text{X (a coloured species in solution)}$\\
$\text{X}\xrightarrow[\text{(2)Pb(CH}_3\text{COO)}_2]{\text{(1) basic medium}}\text{yellow ppt.}$\\
$\text{Q}\xrightarrow[\text{warm}]{\text{MnO(OH)}_2,\ \text{conc. H}_2\text{SO}_4}\text{Y}$ (Y is gas in elemental state and gives blue-coloration with KI-starch paper)\\
A = Number of species among the following which on reaction , changes the color of X aq. solution in acidic medium.\\
H$_2$S , Fe$_2$(SO$_4$)$_3$ , H$_2$C$_2$O$_4$ , FeSO$_4$ , H$_2$O$_2$/H$^+$ , KI\\
B = Number of halides among F$^-$ , Br$^-$ , I$^-$ which can be oxidized by Y\\
Find the sum of (A + B)

\partsec{3}{MATHEMATICS}
\sect{Section-I(i)}

\qn{1} Let $S=\dfrac{1^2}{1.3}+\dfrac{2^2}{3.5}+\dfrac{3^2}{5.7}+\dots+\dfrac{500^2}{999.1001}$, then $[S]$ is less than (where [.] denotes greatest integer function)
\begin{qoptions}
\opt{A}{124}
\opt{B}{126}
\opt{C}{128}
\opt{D}{130}
\end{qoptions}

\qn{2} If $|\vec u|=|\vec v|=1$ and $\vec u.\vec v=0$. Also $[\vec u\ \vec v\ \vec w]=\dfrac54$, $\vec w.\vec u=-\dfrac23\vec w.\vec v$ and $\displaystyle\int_{-2\vec w.\vec u}^{-\frac43\vec w.\vec v}\dfrac{x^3+1}{x^2+1}dx=\dfrac\pi2$. Which of the following is (are) correct ?
\begin{qoptions}
\opt{A}{$|\vec w|=\dfrac74$}
\opt{B}{$\vec w.\vec w=\dfrac{19}8$}
\opt{C}{$\vec w=\dfrac{\vec u}2-\dfrac{3\vec v}2+\dfrac54(\vec u\times\vec v)$}
\opt{D}{$\vec w=\dfrac{\vec u}2-\dfrac{3\vec v}4+\dfrac54(\vec u\times\vec v)$}
\end{qoptions}

\qn{3} Given A and B are two $3\times3$ matrices, such that $A=\text{adj}(B)-B^T$ and $B=\text{adj}A-A^T$. Also, A is a non-singular matrix. Let $|A|+|B|=\alpha$, $AB+BA=\beta I$, then (where I is an Identity matrix of order 3)
\begin{qoptions}
\opt{A}{$\alpha=8$}
\opt{B}{$\alpha=16$}
\opt{C}{$\beta=2$}
\opt{D}{$\beta=8$}
\end{qoptions}

\qn{4} Let $Z_1,Z_2,Z_3,\dots\dots\dots Z_9$ be the distinct roots of $Z^9-9=0$. Let $\lambda=\displaystyle\prod_{1\leqslant i<j\leqslant9}(z_i+z_j)$. Then $\dfrac{\lambda^2}{9^7}$ is
\begin{qoptions}
\opt{A}{3}
\opt{B}{9}
\opt{C}{$9^9$}
\opt{D}{$\dfrac19$}
\end{qoptions}

\qn{5} Let $A=\displaystyle\int_0^\infty e^{-\left(z^2+\frac1{z^2}\right)}dz$, $B=\displaystyle\int_0^\infty e^{-y^2}dy$ then $\dfrac AB$ is not equal to
\begin{qoptions}
\opt{A}{$2e^2$}
\opt{B}{$e^2$}
\opt{C}{$\dfrac1{e^2}$}
\opt{D}{$\dfrac1{2e^2}$}
\end{qoptions}

\qn{6} If $f(x)=\begin{cases}\cos^{-1}\dfrac{1+x}{\sqrt{2(1+x^2)}},&x\leqslant0\\[2mm] \tan^{-1}x,&x>0\end{cases}$\\
If $f(x)=k$ have exactly two real solutions where k is some real number, then
\begin{qoptions}
\opt{A}{$k\in\left(\dfrac\pi2,\pi\right]$}
\opt{B}{k can be $\dfrac\pi3$}
\opt{C}{$k\in\left[\dfrac\pi4,\dfrac\pi2\right)$}
\opt{D}{k can be `1'}
\end{qoptions}

\sect{Section-I(ii)}

\qn{7}\par\nobreak
\begin{center}\small\setlength{\tabcolsep}{3pt}
\begin{tabular}{|c|p{3.9cm}|c|p{1.0cm}|}
\hline
 & \multicolumn{1}{c|}{\textbf{List-I}} & & \textbf{\footnotesize List-II} \\
\hline
(I) & Let $f(x)=\begin{cases}x^{3/5}&\text{if }x\leqslant1\\ -(x-2)^3&\text{if }x>1\end{cases}$\newline Then the number of critical point(s) on the graph of the function is(are) & (P) & 5 \\
\hline
(II) & Number of real solution of the equation, $\log_2{}^2x+(x-1)\log_2x=6-2x$, is & (Q) & 4 \\
\hline
(III) & The number of value(s) of c such that the straight line $3x+4y=c$ touches the curve $\dfrac{x^4}2=x+y$, is(are) & (R) & 3 \\
\hline
(IV) & If $f(x)=\displaystyle\int_x^{x^2}(t-1)\,dt$, $1\leqslant x\leqslant2$, then global maximum value of $f(x)$ is & (S) & 2 \\
\hline
 & & (T) & 1 \\
\hline
\end{tabular}
\end{center}
\begin{qoptions}
\opt{A}{I $\to$ R; II $\to$ S; III $\to$ T; IV $\to$ Q}
\opt{B}{I $\to$ R; II $\to$ T; III $\to$ S; IV $\to$ P}
\opt{C}{I $\to$ R; II $\to$ Q; III $\to$ P; IV $\to$ S}
\opt{D}{I $\to$ R; II $\to$ Q; III $\to$ T; IV $\to$ S}
\end{qoptions}

\qn{8} Consider the quadratic trinomial $f(x)=2x^2-10px+7p-1$, where p is a parameter. Match the range of p in List-II with the conditions given in List-I.\par\nobreak
\begin{center}\small\setlength{\tabcolsep}{3pt}
\begin{tabular}{|c|p{3.6cm}|c|p{2.3cm}|}
\hline
\multicolumn{2}{|c|}{\textbf{List-I}} & \multicolumn{2}{c|}{\textbf{List-II}} \\
\hline
(I) & If both roots of $f(x)=0$ are confined in $(-1,1)$ then & (P) & $\left(\dfrac25,\infty\right)$ \\
\hline
(II) & Exactly one root of $f(x)=0$ lies in $(-1,1)$ & (Q) & $\phi$ \\
\hline
(III) & Both roots of $f(x)=0$ are greater than I & (R) & $\left(-\dfrac1{17},\dfrac13\right)$ \\
\hline
(IV) & One root of $f(x)=0$ is greater than 1 and other root of $f(x)=0$ is less than $-1$ & (S) & $\left(-\infty,-\dfrac1{17}\right]\cup\left[\dfrac13,\infty\right)$ \\
\hline
 & & (T) & $\left(-\infty,-\dfrac1{17}\right)\cup\left(\dfrac13,\infty\right)$ \\
\hline
\end{tabular}
\end{center}
\begin{qoptions}
\opt{A}{I $\to$ R; II $\to$ Q; III $\to$ S; IV $\to$ T}
\opt{B}{I $\to$ R; II $\to$ S; III $\to$ Q; IV $\to$ Q}
\opt{C}{I $\to$ R; II $\to$ T; III $\to$ S; IV $\to$ Q}
\opt{D}{I $\to$ R; II $\to$ P; III $\to$ Q; IV $\to$ T}
\end{qoptions}

\qn{9} Match the entries of List-I with one or more than one entries of List-II.\par\nobreak
\begin{center}\small\setlength{\tabcolsep}{3pt}
\begin{tabular}{|c|p{4.5cm}|c|p{1.2cm}|}
\hline
\multicolumn{2}{|c|}{\textbf{List-I}} & \multicolumn{2}{c|}{\textbf{List-II}} \\
\hline
(I) & Let $f:[-1,1]\to R$ be defined by $f(x)=\sqrt[5]x+\sin^{-1}x$ then $f(x)$ is & (P) & Odd \\
\hline
(II) & Let $f:R\to\{-1,0,1\}$ is defined by $f(x)=\text{sgn}\left(\dfrac{1-|x|}{1+|x|}\right)$ then $f(x)$ is (sgn(x) denotes signum function of x) & (Q) & Even \\
\hline
(III) & Let $f:(-\infty,0]\to[0,\infty)$ be defined by $f(x)=\dfrac{2^{-[x]}}{2^{\{x\}}}-2^{|x|}$, then $f(x)$ is ($[x]$ denotes the largest integer less than or equal to x and $\{x\}$ denotes fractional part of x) & (R) & Onto \\
\hline
(IV) & Let $f:(-4,2]\to[0,3]$ be defined by $f(x)=\sqrt{8-2x-x^2}$, then $f(x)$ is & (S) & One-one \\
\hline
 & & (T) & Many-one \\
\hline
\end{tabular}
\end{center}
\begin{qoptions}
\opt{A}{I $\to$ P,S; II $\to$ Q,R,T; III $\to$ R,T; IV $\to$ R,T}
\opt{B}{I $\to$ P,S; II $\to$ R,T; III $\to$ Q,R,T; IV $\to$ T}
\opt{C}{I $\to$ P,S; II $\to$ R,T; III $\to$ T; IV $\to$ Q,R,T}
\opt{D}{I $\to$ P,S; II $\to$ T; III $\to$ Q,R,T; IV $\to$ R,T}
\end{qoptions}

\qn{10}\par\nobreak
\begin{center}\small\setlength{\tabcolsep}{3pt}
\begin{tabular}{|c|p{4.0cm}|c|p{0.7cm}|}
\hline
\multicolumn{2}{|c|}{\textbf{List-I}} & \multicolumn{2}{c|}{\textbf{List-II}} \\
\hline
(I) & If the mid point of a chord of the ellipse $\dfrac{x^2}9+\dfrac{y^2}{16}=1$ is $(0,3)$ then length of the chord is $\dfrac ab\sqrt c$, then the value of $(a+b+c)$ (gcd (a, b) = 1, gcd (b, c) = 1, gcd (a, c) = 1 where gcd denotes greatest common divisor) & (P) & 1 \\
\hline
(II) & PQ is a chord of the parabola $y^2=4x$ joining $P(1,2)$ \& $Q(t^2,2t)$ PQ subtend $90^\circ$ at origin then $|t|$, is & (Q) & 12 \\
\hline
(III) & The number of circle's touching all the three lines $3x+8y=7$, $21x+56y=9$, $6x+16y=0$ & (R) & 4 \\
\hline
(IV) & The axes are translated so that the new equation of the circle $x^2+y^2-7x+9y+10=0$ has no first degree term \& the new equation is $x^2+y^2=\dfrac\lambda2$ then the value of $\lambda$ is & (S) & 45 \\
\hline
 & & (T) & 0 \\
\hline
\end{tabular}
\end{center}
\begin{qoptions}
\opt{A}{I $\to$ Q; II $\to$ R; III $\to$ T; IV $\to$ S}
\opt{B}{I $\to$ Q; II $\to$ R; III $\to$ R; IV $\to$ S}
\opt{C}{I $\to$ P; II $\to$ R; III $\to$ Q; IV $\to$ S}
\opt{D}{I $\to$ P; II $\to$ R; III $\to$ Q; IV $\to$ T}
\end{qoptions}

\sect{Section-II}

\qn{1} Let $\displaystyle\lim_{x\to\infty}x\ln\left(e\left(1+\dfrac1x\right)^{1-x}\right)$ equals $\dfrac mn$ where m and n are relatively prime positive integer. Find $2m-n$.

\qn{2} A man has to distribute 8 different chocolates among his three grand children such that no two of grand children receive same number of chocolates but each receives atleast one chocolate and none is left over, then the number of ways in which the division can be made ?

\qn{3} From a lot containing 4 defective and 6 good items, 4 items are chosen at random which are to be checked for their quality. As a result of this inspection, first 2 of the 4 items chosen found to be both good then the probability that 4 chosen items for inspection were all good, can be expressed as $\dfrac pq$ where p, q are coprime, then the value of $(p+q)$ is

\qn{4} If $|z|=\min\{|z-2|,|z+2|\}$, then the value of $|z+\bar z|$, is

\qn{5} Let point $A(x_1,y_1,0)$ and $B(x_2,y_2,0)$ are lying on curve, $y=2+e^{\ln(1-x)}+\log_{1/2}\left(2^{\sqrt{x^2-2x+1}}\right)$ and A also lies on curve $x^2+y^2=13$, $z=0$ and B lies inside it. If $x_1,y_1,x_2,y_2$ are integers and for largest value of $x_2$ shortest distance of A from line $\dfrac{x-x_1}{x_2+3}=\dfrac{y-y_2}4=\dfrac{z-3}5$ is $\dfrac a{\sqrt b}$, a and b are coprime, then the value of $(a+b)$, is

\qn{6} Let $I(a,b)=\displaystyle\int\dfrac{x^b}{(1+x^2)^a}dx$\\
$I(a,b)=A\dfrac{x^{b-1}}{(1+x^2)^{a-1}}+B\,I(a,b-2)$\\
Then the value of $(A+B)$, is (where a = 50, b = 97)

\qn{7} Let, tangent to a curve $y=f(x)$ at some random point P, intersect the y-axis at A and normal to the same curve at same point P, intersect x-axis at B if circumcenter of triangle PAB lies on the line $y=x$. If $y=f(x)$ satisfies the equation, $\ln\left(x^2+y^2\right)^{\frac1n}+\tan^{-1}\dfrac yx=k$ with $f(1)=0$. Then value of $(n+k)$ is

\qn{8} Let $g(k,z)=Z^k\displaystyle\int_0^{\frac1z}t^{t+k-1}\,dt$. If $\displaystyle\lim_{z\to\infty}g(k,z)=f(k)$, then number of point(s) of intersection of curves $xy=\dfrac1{f(4)}$ and $x^2+y^2=8$ is(are)

\papersec{2}
\partsec{1}{PHYSICS}
\sect{Section-I(i)}

\qn{1} Figure shows a transparent block of front sides `a' and `a'. The third dimension of the block is negligible. A point source S which can emit light in all directions can move inside the block. It is desired that no light of `S' should pass through AB. The region in which S should be present to satisfy this condition shown by shaded region. Choose correct option.
\qfig{p2_phy_s1_q01_1}
\begin{qoptions}
\opt{A}{\qfigi{p2_phy_s1_q01_2}}
\opt{B}{\qfigi{p2_phy_s1_q01_B}}
\opt{C}{\qfigi{p2_phy_s1_q01_C}}
\opt{D}{\qfigi{p2_phy_s1_q01_D}}
\end{qoptions}

\qn{2} A string on a musical instrument is held under tension T and extends from the point x = 0 to point x = L. The string is overwound with a wire in such a way that its mass per unit length $\mu$ (x) increases uniformly from $\mu_0$ at x = 0 to $\mu_L$ at x = L. Find the time interval required for a transverse pulse to travel the length of the string.
\begin{qoptions}
\opt{A}{$\Delta t=\dfrac{3L\left(\mu_L+\mu_0+\sqrt{\mu_L\mu_0}\right)}{2\sqrt T\left(\sqrt{\mu_L}-\sqrt{\mu_0}\right)}$}
\opt{B}{$\Delta t=\dfrac{3L\left(\mu_L+\mu_0-\sqrt{\mu_L\mu_0}\right)}{2\sqrt T\left(\sqrt{\mu_L}+\sqrt{\mu_0}\right)}$}
\opt{C}{$\Delta t=\dfrac{3L\left(\mu_L+\mu_0+\sqrt{\mu_L\mu_0}\right)}{2\sqrt T\left(\sqrt{\mu_L}+\sqrt{\mu_0}\right)}$}
\opt{D}{$\Delta t=\dfrac{2L\left(\mu_L+\mu_0+\sqrt{\mu_L\mu_0}\right)}{3\sqrt T\left(\sqrt{\mu_L}+\sqrt{\mu_0}\right)}$}
\end{qoptions}

\qn{3} A uranium nucleus ${}^{235}$U liberates 200 MeV per fission. 1.5 kg of uranium reacts during explosion of a uranium bomb. What is the mass of an equivalent TNT bomb if the heating capacity of TNT is $4.1\times10^6$ J/kg?
\begin{qoptions}
\opt{A}{$3\times10^6$ kg}
\opt{B}{$3\times10^7$ kg}
\opt{C}{$9\times10^7$ kg}
\opt{D}{$9\times10^6$ kg}
\end{qoptions}

\qn{4} A solid sphere (density = 8900 kg/m$^3$ and specific heat capacity S = 390 J-kg$^{-1}$-K$^{-1}$) of radius 10 cm is at an initial temperature $T_1$ = 200 K. It is then suspended inside a chamber whose walls are at almost 0 K. Calculate the time required for the temperature of the sphere to drop to $T_2$ = 100 K. Given $\sigma=5.67\times10^{-8}\,\text{W/m}^2-\text{K}^4$
\begin{qoptions}
\opt{A}{123 hrs 21 mins}
\opt{B}{227 hrs 51 mins}
\opt{C}{148 hrs 23 mins}
\opt{D}{165 hrs 19 mins}
\end{qoptions}

\sect{Section-I(ii)}

\qn{5} Total energy of a particle executing oscillating motion is 3 joule and given by\\
E = $x^2+2x+v^2-2v$\\
Where x is the displacement from origin and v is velocity of particle at x. Then choose the correct statement(s) :
\begin{qoptions}
\opt{A}{Amplitude of oscillation is 1m}
\opt{B}{Maximum velocity of particle during oscillation is 3 m/sec}
\opt{C}{Amplitude of oscillation can be 2m}
\opt{D}{Maximum velocity of particle during oscillation is $\left(\sqrt5+1\right)$ m/sec}
\end{qoptions}

\qn{6} Diagram shows a potentiometer. The wire of potentiometer is of length $\ell$ and is of uniform cross-section A but its resistivity varies as $\rho=\rho_0\dfrac x\ell$ (x measured from end P ). When switch was opened, balancing length was $\ell_1$ while when switch was closed, balancing length was $\ell_2$. Choose the correct option(s) :
\qfig{p2_phy_s2_q06_1}
\begin{qoptions}
\opt{A}{Value of $\varepsilon_1$ is $\dfrac{\rho_0\varepsilon_0\ell_1^2}{2A\ell r_0+\rho_0\ell^2}$}
\opt{B}{Value of $\varepsilon_1$ is $\dfrac{\rho_0\varepsilon_0\ell_1\ell_2}{\frac R{\ell_2}+2A\ell r_0+\rho_0\ell_1\ell_2}$}
\opt{C}{Value of r is $R\left[\left(\dfrac{\ell_1}{\ell_2}\right)^2-1\right]$}
\opt{D}{Value of r is $R\left[\left(\dfrac{\ell_2}{\ell_1}\right)^2-1\right]$}
\end{qoptions}

\qn{7} The figure shows a cylindrical container fitted with a piston. It contains two immiscible incompressible liquids with densities $\rho_1$ and $\rho_2$. A block floats is equilibrium in the position shown. Choose the \textbf{CORRECT} statement(s). (Spring is connected to bottom of cylinder as well as block.)
\qfig{p2_phy_s2_q07_1}
\begin{qoptions}
\opt{A}{Suppose the spring is compressed. Now, if the piston is moved downward then compression in spring will be same as before.}
\opt{B}{Suppose the spring is compressed. Now if the whole system is moved with some acceleration upward then spring will definitely attain natural length.}
\opt{C}{Supposed the spring is initially relaxed. The piston is moved downward and whole system is accelerated upward then no change will occur in the state of spring.}
\opt{D}{Supposed the spring is initially relaxed. The piston is moved downward and whole system is accelerated upward then the change in state of spring will depend on $\rho_1$ and $\rho_2$.}
\end{qoptions}

\qn{8} Consider the PV cycle shown. If the heat given to one mole of the monatomic gas in the step AB is Q, then :
\qfig{p2_phy_s2_q08_1}
\begin{qoptions}
\opt{A}{Work done in the process BC on the gas is 3Q /10}
\opt{B}{The heat released in the process BC is 2Q/5}
\opt{C}{The efficiency of the cycle is 10\%}
\opt{D}{The heat given to the gas during CA is 3Q/10}
\end{qoptions}

\qn{9} A particle having charge Q and mass M = 2m is tied to two identical particles, each having mass m and charge q. The particles are connected with ideal strings of equal length $\ell$. The system is held at rest on a smooth horizontal surface (with string taut) in a position where string make an angle of $60^\circ$ between them. From this position the system is released :
\qfig{p2_phy_s2_q09_1}
\begin{qoptions}
\opt{A}{The amplitude of oscillation of M is $\dfrac{\sqrt3}4\ell$}
\opt{B}{Maximum speed acquired by M is $\dfrac q4\sqrt{\dfrac1{\pi\epsilon_0m\ell}}$}
\opt{C}{When all the particles are in straight line, speed of each particle is same.}
\opt{D}{Tension in string when particles are in straight line is $T=\dfrac q{16\pi\epsilon_0\ell^2}(4Q+5q)$}
\end{qoptions}

\qn{10} A dumb-bell is made of two small spheres of equal mass connected by a light rigid rod of length $2\ell$. Two such identical dumb-bells approach each other with zero angular velocity on a smooth horizontal table as shown and undergo an elastic collision of negligible duration. Which of the following statements is/are true?
\qfig{p2_phy_s2_q10_1}
\begin{qoptions}
\opt{A}{Just after collision, the tension in the rods is $\dfrac{2mv^2}\ell$}
\opt{B}{The rods do not exert any impulse on the spheres during the collision}
\opt{C}{The center of mass of each dumb-bell remains at rest for a time $\dfrac{\pi\ell}v$}
\opt{D}{There will be a second collision.}
\end{qoptions}

\sect{Section-III}

\qn{1} A parallel plate capacitor with plates of length $\ell$ is included in a circuit as shown in figure. Given are the emf of the current source, its internal resistance r and the distance d between the plates. An electron with a velocity $V_0$ flies into the capacitor, parallel to the plates. The resistance R should be connected in parallel with the capacitor so that the electron flies out of the capacitor at an angle of $37^\circ$ to the plates is given as $\alpha\ \Omega$. Find $\dfrac\alpha{10}$.\\
(Assume that circuit is in steady state). (Take : $\ell$ = 91 cm, $\varepsilon$ = 3V, r = $2\Omega$, d = $\dfrac13$mm, $v_0=4\times10^7$ m/s)
\qfig{p2_phy_s3_q01_1}

\qn{2} A point$-$like light$-$source is placed (in air) to the symmetry axis of a hemisphere shaped piece of glass of radius 20 cm, 10 cm away from the plane surface. The refractive index of glass is 1.5. If a ray travels a distance x (in cm) through the glass till it reaches the curved surface as shown in figure then find the value of $\dfrac{\sqrt2+1}5$ x.
\qfig{p2_phy_s3_q02_1}

\qn{3} Two point charges A and B, having charges $q_A$ and $q_B$ in ratio $\dfrac{q_A}{q_B}=\dfrac{-\left(\sqrt2+1\right)}{\sqrt2-1}$, enter a uniform magnetic field given by $\vec B=0$ for x $<$ 0 \& $\vec B=-B_0\hat k$ for x $>$ 0. Charges enter at origin making an angle 45$^\circ$ with x-axis moving in xy plane as shown, one by one, and eventually, come out. During their motion the maximum x-coordinate acquired by both of them is equal. If their de-Broglie wavelengths are $\lambda_A$ and $\lambda_B$ respectively, find $\dfrac{\lambda_A}{\lambda_B}$. (Ignore interaction between A and B) ($q_A$ is $-$ve)
\qfig{p2_phy_s3_q03_1}

\qn{4} Consider a standard YDSE setup in which a bichromatic source of light is used. Find the 2nd closest (non-zero) distance (in mm) from the central bright where the intensity is maximum. (the terms have their usual meaning)\\
(Given : $\lambda_1$ = 300 nm, $\lambda_2$ = 500 nm, d = 1mm, D = 1m)

\qn{5} There exists a uniform magnetic field perpendicular to the plane of the figure in cylindrical region of radius $a$. The magnetic field is increasing at a constant rate $\alpha\,\text{Ts}^{-1}$. A particle having charge q is at a point $P$ outside the field region. The particle is slowly moved to infinity, work done by the external agent on the particle if, the particle is moved in a direction perpendicular to $OP$ along $PQ$ is $\dfrac{\pi qa^2\alpha}n$. Find n.
\qfig{p2_phy_s3_q05_1}

\qn{6} In a hypothetical hydrogen atom the electrostatic potential energy of interaction of proton and electron is given by $U=U_0\ln\left(\dfrac r{r_0}\right)$ where $U_0$ and $r_0$ are constants and r is radius of circular orbit of electron. For such hydrogen atom the energy difference between $n^{\text{th}}$ and $m^{\text{th}}$ state is represented by $\Delta E_{nm}$. Calculate the ratio $\Delta E_{12}:\Delta E_{24}$. Assume Bohr's assumption of angular momentum quantization to hold.

\qn{7} In a region of space a uniform magnetic field exist in positive z direction and there also exists a uniform electric field along positive y direction. A particle having charge +q and mass m is released from rest at the origin. The particle moves on a curve known as cycloid. If a wheel of radius R were to roll on the X axis, a fixed point on the circumference of the wheel would generate this cycloid. It is given that strength of magnetic and electric fields are $B_0$ and $E_0$ respectively. Radius, R = $\dfrac{nmE_0}{7qB_0^2}$. Find n.

\qn{8} A solid sphere of mass m and radius r is given an initial angular velocity $\omega_0$ and a linear velocity $v_0=\lambda r\omega_0$ from a point A on a rough horizontal surface. It is observed that the ball turns back and returns to the point A after some time if $\lambda$ is less than a certain maximum value $\lambda_0$. Find $10\lambda_0$.
\qfig{p2_phy_s3_q08_1}

\partsec{2}{CHEMISTRY}
\sect{Section-I(i)}

\qn{1} When 0.1 mole salt of cobalt (III) forming octahedral structure with chemical formula $\left[\text{Co(NH}_3)_x\text{Cl}_y\right]\text{Cl}_z$ get dissolved in 31 mole of water, then freezing point of solution becomes 272.15K. Determine the value of (x + y + z), if $K_f$ of water is 1.86 K kg mol$^{-1}$ and freezing point of pure water is 273.15 K (assuming dissociation of given salt is 100\%)
\begin{qoptions}
\opt{A}{6}
\opt{B}{7}
\opt{C}{8}
\opt{D}{9}
\end{qoptions}

\qn{2} Which of the following reaction will produce HD ? $\left(\text{D}={}^2_1\text{H}\right)$ without contamination of H$_2$ and D$_2$?
\begin{qoptions}
\opt{A}{Applying electric arc to a mixture of H$_2$ \& D$_2$}
\opt{B}{Adding D$_2$O to NaH}
\opt{C}{Electrolysis of HDO}
\opt{D}{Electrolysis of NaCl solution in D$_2$O.}
\end{qoptions}

\qn{3} The correct statement among the following is :-
\begin{qoptions}
\opt{A}{$\text{WO}_4^{2-}>\text{MoO}_4^{2-}>\text{CrO}_4^{2-}$ (oxidising power)}
\opt{B}{Ni = Pt $>$ Pd (No. of unpaired `d' electrons)}
\opt{C}{Fe $>$ Co $>$ Ni (density)}
\opt{D}{Eu $>$ Gd $>$ Tb (No. of protons)}
\end{qoptions}

\qfullbegin
\qn{4} \qfigi{p2_chem_s1_q04_1}\\
Major product `\textbf{E}' is
\begin{qoptions}
\opt{A}{\qfigi{p2_chem_s1_q04_3}}
\opt{B}{\qfigi{p2_chem_s1_q04_B}}
\opt{C}{\qfigi{p2_chem_s1_q04_C}}
\opt{D}{\qfigi{p2_chem_s1_q04_D}}
\end{qoptions}
\qfullend

\sect{Section-I(ii)}

\qn{5} The standard boiling point of benzene (C$_6$H$_6$) is 80.1$^\circ$C, then for given process\\
C$_6$H$_6$ ($\ell$) $\rightarrow$ C$_6$H$_6$ (g)\\
Select the correct statement(s):
\qfig{p2_chem_s2_q05_1}
\begin{qoptions}
\opt{A}{$\Delta$G is zero at standard boiling point for the given process}
\opt{B}{$\Delta$G$^\circ$ is zero at standard boiling point for the given process}
\opt{C}{$\Delta$H$^\circ$ is positive for given process}
\opt{D}{$\Delta$S$^\circ$ is positive for given process.}
\end{qoptions}

\qn{6} For the gas phase reaction :\\
C$_2$H$_4$ + H$_2$ $\rightleftharpoons$ C$_2$H$_6$ $(\Delta\text{H}=-32.7\,\text{Kcal})$\\
carried out in a vessel, the equilibrium concentration of C$_2$H$_4$ can be increased by-
\begin{qoptions}
\opt{A}{increasing the temperature}
\opt{B}{decreasing the pressure}
\opt{C}{removing some H$_2$ at constant volume}
\opt{D}{adding some C$_2$H$_6$ at constant volume}
\end{qoptions}

\qn{7} Chiral compounds among the following are :-
\begin{qoptions}
\opt{A}{$\left[\text{Ru(bpy)}_3\right]^{2+}$}
\opt{B}{Fac $-$ $\left[\text{Co(NO}_2)_3\text{ (dien)}\right]$}
\opt{C}{Mer $-$ $\left[\text{Co(NO}_2)_3\text{ (dien)}\right]$}
\opt{D}{$\left[\text{Mn(acac)}_3\right]$}
\end{qoptions}

\qn{8} Which of the following is correctly matched ?
\begin{qoptions}
\opt{A}{Zone refining - Ga, In}
\opt{B}{Vapour phase refining - Ti, Ni}
\opt{C}{Liquation - Sn, W}
\opt{D}{Distillation - Zn, Hg}
\end{qoptions}

\qn{9}
\qfig{p2_chem_s2_q09_1}
Select the correct options.
\begin{qoptions}
\opt{A}{Product `C' is a symmetrical diketone.}
\opt{B}{`D' produces brisk effervescence with NaHCO$_3$.}
\opt{C}{Formation of `A' is an example of conjugate addition.}
\opt{D}{`D' on monobromination with Br$_2$/h$\upsilon$ produces 3 products (not including stereoisomers)}
\end{qoptions}

\qfullbegin
\qn{10} Which of the following reactions are feasible?
\begin{qoptions}
\opt{A}{\qfigi{p2_chem_s2_q10_1}}
\opt{B}{\qfigi{p2_chem_s2_q10_B}}
\opt{C}{\qfigi{p2_chem_s2_q10_C}}
\opt{D}{\qfigi{p2_chem_s2_q10_D}}
\end{qoptions}
\qfullend

\sect{Section-III}

\qn{1} At 27$^\circ$C and 1 bar, gaseous N$_2$O$_4$ is 50\% dissociated into NO$_2$ at equilibrium by the reaction :\\
N$_2$O$_4$ $\rightleftharpoons$ 2NO$_2$\\
If magnitude of standard Gibbs free energy change of the reaction at 727$^\circ$C is `P' cal/mol, then find value of $\left(\dfrac P{100}\right)$ ?\\
\relax[Given : ln2 = 0.7, ln3 = 1.1]

\qn{2} The standard reduction potential of silver chloride electrode (Ag/AgCl/Cl$^-$)(metal $-$ sparingly soluble salt electrode) is 0.21 V and for silver electrode is 0.80 V. If the moles of AgCl that can be dissolved in 10 L of a 0.01 M NaCl solution is represented as $10^{-y}$, then find the value of y?\\
$\left[\text{Given : }\dfrac{2.303\,RT}F=0.059\right]$

\qn{3} For the reaction\\
Br$_2$($l$) + Cl$_2$(g) $\rightleftharpoons$ 2BrCl(g) ;\\
In a closed container of volume 164 L initially 10 moles of Cl$_2$ are present at 27$^\circ$C. What minimum mass (in gram) of Br$_2$($l$) must be introduced into this container so that above equilibrium is maintained at total pressure of 2.25 atm. Vapour pressure of Br$_2$($l$) at 27$^\circ$C is 0.25 atm. Assume that volume occupied by liquid is negligible. (R = 0.082 L atm mol$^{-1}$ K$^{-1}$, Atomic mass of Br = 80). Express minimum mass required as $x\times10^2$ g. Determine $x$.

\qn{4} Pyrolusite, Rutile, Chromite, Bauxite, Kernite, Limonite, Chalcopyrite, Cerrusite.\\
Total number of transition metal oxide ores are :-

\qn{5} The oxidation state of Nickel in the final product of the given reaction is :-\\
$\text{Ni}^{+2}\text{(aq.)}\xrightarrow{\text{KCN}}\text{A}\xrightarrow{\text{K metal in liq. NH}_3}\text{B}$

\qn{6}
\qfig{p2_chem_s3_q06_1}
Final product (D) is soluble in aq. NaHCO$_3$. Calculate the total number of sp$^2$ hybridised carbons in (D).

\qn{7} Find total number of reactions in which final product can be secondary amine.
\qfig{p2_chem_s3_q07_1}

\qn{8}
\qfig{p2_chem_s3_q08_1}
How many given statements are correct regarding final product?\\
(1) At pH=0, it is colorless and total charge of this colorless compound is 0.\\
(2) At pH=14, it is colorless and total charge of this colorless compound is $-3$.\\
(3) Number of sp$^2$ hybridised carbon atoms in final product is 20.\\
(4) Number of acidic hydrogen atoms in final product is 3.

\partsec{3}{MATHEMATICS}
\sect{Section-I(i)}

\qn{1} For $\theta\in(0,2\pi)-\{\pi\}$, sum of all values of $`\theta'$ satisfying equation $4\cot^4\theta+9\cot^3\theta+8\cot^2\theta+9\cot\theta+4=0$; is equal to :
\begin{qoptions}
\opt{A}{$5\pi$}
\opt{B}{$6\pi$}
\opt{C}{$7\pi$}
\opt{D}{None}
\end{qoptions}

\qn{2} For p, q, r $>$ 0 the least value of $\dfrac{\left(2p^8+2q^8+3r^2\right)^2}{p^2q^2r^3}$ is equal to:
\begin{qoptions}
\opt{A}{16}
\opt{B}{$16\sqrt2$}
\opt{C}{32}
\opt{D}{$32\sqrt2$}
\end{qoptions}

\qn{3} How many of following statement(s) is/are TRUE?\\
(i) There exist infinite pairs of square matrices (A, B) such that $AB-BA=I$, where I is identity matrix of same order.\\
(ii) For two invertible matrices of same order A \& B det(I + AB) = det(I + BA)\\
(iii) Determinant value of each skew-symmetric matrix with integral entries is perfect square of an integer.\\
(iv) For distinct positive integers a, b, c, d if $\begin{vmatrix}a&0&1\\a&b&-b\\0&1&c\end{vmatrix}=23$, then least of a + b + c is 13.
\begin{qoptions}
\opt{A}{1}
\opt{B}{2}
\opt{C}{3}
\opt{D}{4}
\end{qoptions}

\qn{4} Value of $\displaystyle\sum_{j=1}^{10}{}^{10}C_j\left(\sum_{i=1}^{10-j}{}^{10-j}C_i\right)$ is divisible by :
\begin{qoptions}
\opt{A}{5}
\opt{B}{3}
\opt{C}{2}
\opt{D}{None}
\end{qoptions}

\sect{Section-I(ii)}

\qn{5} If every solution of the equation $x([x]-5)+2\{x\}+6=0$, satisfies the inequality $(a-3)x^2+2(a+3)x-8a\leqslant0$, where $[x]$ and $\{x\}$ denote greatest integer and fractional part functions respectively then value(s) of a can be
\begin{qoptions}
\opt{A}{$(-2,1)$}
\opt{B}{$(-9,0]$}
\opt{C}{$\{0,1\}$}
\opt{D}{$(1,2)$}
\end{qoptions}

\qn{6} If\\ \resizebox{\linewidth}{!}{$f(x)=\displaystyle\lim_{n\to\infty}\dfrac{\left(1-\cos\left\{1-\tan\left(\frac\pi4-x\right)\right\}\right)(x+1)^n+\lambda\sin\left(n-\sqrt{n^2-8n}\right)x}{x^2(x+1)^n+x}$}\\
$x\neq0$ is continuous at $x=0$ then
\begin{qoptions}
\opt{A}{$f(0)=1$}
\opt{B}{$\lambda=\dfrac12$}
\opt{C}{$f(0)=2$}
\opt{D}{$\lambda=1$}
\end{qoptions}

\qn{7} Consider the family of planes $x+y+z=c$ where c is a parameter intersecting the coordinate axes at P, Q, R and $\alpha,\beta,\gamma$ are the angles made by each member of this family with positive x, y and z axis. Which of the following interpretations hold good for this family.
\begin{qoptions}
\opt{A}{each member of this family is equally inclined with the coordinate axes.}
\opt{B}{$\sin^2\alpha+\sin^2\beta+\sin^2\gamma=1$}
\opt{C}{$\cos^2\alpha+\cos^2\beta+\cos^2\gamma=2$}
\opt{D}{for c = 3 area of the triangle PQR is $3\sqrt3$ sq. units.}
\end{qoptions}

\qn{8} If $\displaystyle\int\dfrac{(x-1)dx}{x^2\sqrt{2x^2-2x+1}}=\dfrac{\sqrt{f(x)}}{g(x)}+c$, then
\begin{qoptions}
\opt{A}{$f(x)=2x^2-2x+1$}
\opt{B}{$g(x)=x+1$}
\opt{C}{$f(x)=2x^2-2x$}
\opt{D}{$g(x)=x$}
\end{qoptions}

\qn{9} Let $f(x,y)=0$ is a solution of the differential equation $\dfrac{dy}{dx}=\dfrac{y-2x^3}{x+2yx^2}$ in x-y plane where $f(1,1)=3$. The value of $\lambda$, for which the line $\dfrac{x-1}2=\dfrac{y+1}3=\dfrac{z+1}1$ intersects the curve $f(x,y)=\lambda$, is less than :
\begin{qoptions}
\opt{A}{7}
\opt{B}{14}
\opt{C}{13}
\opt{D}{22}
\end{qoptions}

\qn{10} Let $f:R\to R$ be a differentiable function satisfying $f\left(\dfrac{x+y}2\right)=\dfrac{f(x)+f(y)}2\ \forall\ x,y\in R$. If $f(0)=1$ and $f'(0)=-1$, then which of the following is/are correct ?
\begin{qoptions}
\opt{A}{$f(|x|)$ is non-differentiable at one point}
\opt{B}{Number of solution of the equation $f(x)=f^{-1}(x)$ is exactly 1.}
\opt{C}{$\displaystyle\sum_{r=0}^{10}(f(r))^2=286$}
\opt{D}{$\tan^{-1}(f(x))$ is derivable $\forall\ x\in R$}
\end{qoptions}

\sect{Section-III}

\qn{1} The area of the loop of the curve $3y^2=x^2(3-x)$ is `A' then find [A] where [ ] represent G.I.F.

\qn{2} A wire of given length `$\ell$' is cut into two pieces which are then bent into two shapes i.e. circle \& square respectively. If the sum of the area of circle and square will be least when (side of square) = k (diameter of circle), then find the sum of digits of $\sqrt{2025}k$.

\qn{3} Value of $\displaystyle\lim_{x\to0}\dfrac{\sin^{-1}x-\tan^{-1}x+\left(1+x^3\right)^{1/3}-1}{x^3}$ is K then value of 6K is

\qn{4} If a, b, c, d $>$ 0, x $\in$ R and $\left(a^2+b^2+c^2\right)x^2-2(ab+bc+cd)x+b^2+c^2+d^2\leqslant0$, then value of $\begin{vmatrix}33&14&\ell na\\65&27&\ell nb\\97&40&\ell nc\end{vmatrix}+\dfrac{b^2+c^2}{ac+bd}$ is

\qn{5} For $\lambda\in Z^+$ (the set of all positive integers),\\
\resizebox{\linewidth}{!}{$\displaystyle\lim_{n\to\infty}\dfrac{(n+7)^{2-\lambda}\left(1^\lambda+2^\lambda+\dots\dots+n^\lambda\right)}{\left(n^2\lambda+1^2\right)+\left(n^2\lambda+2^2\right)+\dots\dots\dots+\left(n^2\lambda+n^2\right)}=\dfrac3{40}$}\\
then $\lambda=$

\qn{6} Let z is a complex number satisfying $\arg\left(\dfrac{z+\ell}{z+m}\right)=\dfrac\pi6$, $(\ell,m\in R)$. Locus of z is an arc of a circle $x^2+y^2+3\sqrt2x-\sqrt6y+4=0$ then $\ell^2+m^2-\ell m=\dots\dots\dots$

\qn{7} The slope of tangent and normal to the curve $\left(y+x^3\right)^3=x^2\left(\tan^{-1}(\tan x)+\text{cosec}^{-1}(\text{cosec }x)+x^2\right)^2$ at the point $(2,-4)$ are $m_1$ and $m_2$ respectively. Then $[m_2-m_1]=\dots\dots\dots$, ([$\cdot$] denotes GIF)

\qn{8} Let S (a, b, c) be the mirror image of a point Q (3, 2, 9) with respect to the plane given by $\vec r=(1+2\mu)\hat i+(2+\lambda-\mu)\hat j-\lambda\hat k$, where $\lambda$ and $\mu$ are parameters. Then c $-$ b = .........

\end{multicols}
\end{document}
